\begin{table*}[t]
\centering
\caption{Melhor configuração escolhida na validação e desempenho no teste. Critério de seleção fixado antes do teste: maior Dice médio de validação entre as 198 combinações.}
\label{tab:best_configuration}
\small
\begin{tabularx}{\textwidth}{L{4.3cm} L{1.6cm} L{2.4cm} L{2.8cm} c c c}
\toprule
\textbf{Conjunto} & \textbf{Arquitetura} & \textbf{Pré-processamento} & \textbf{Pós-processamento} & \textbf{IoU} & \textbf{Dice} & \textbf{IC 95\% (Dice)} \\
\midrule
\texttt{gray\_scott\_reaction\_}\newline\texttt{diffusion} & U-Net & P2 RMS & Q1 Otimização do limiar & 0,950 & 0,962 & [0,960; 0,964] \\
\bottomrule
\end{tabularx}
\vspace{6pt}
\caption{Resumo entre conjuntos da abordagem de segmentação escolhida. Ganho relativo: Dice no teste da configuração escolhida em relação à mesma arquitetura sem pré e pós-processamento.}
\label{tab:cross_dataset}
\small
\begin{tabularx}{\textwidth}{L{4.3cm} L{1.7cm} L{1.4cm} c c Y}
\toprule
\textbf{Conjunto} & \textbf{Arquitetura escolhida} & \textbf{P / Q escolhidos} & \textbf{Dice no teste} & \textbf{Ganho relativo (\%)} & \textbf{Observações de generalização} \\
\midrule
\texttt{gray\_scott\_reaction\_}\newline\texttt{diffusion} & U-Net & P2 / Q1 & 0,962 & +0,95 & O ganho sobre a U-Net sozinha (0,953) vem do instante inicial aleatório; em $n\geq50$, Dice de 0,992 contra 0,995 ($-$0,2\%). \\
\bottomrule
\end{tabularx}
\end{table*}
